%==============================================================================
%  APPENDICES
%  A: Reproducibility and how to run   |   B: AI-assisted development
%==============================================================================
\appendix

\chapter{Reproducibility and How to Run}
\label{ch:reproduce}
The project is self-contained and portable: it targets \textbf{Python~3.11+}, declares
every dependency in \code{pyproject.toml}, and the dashboard runs fully offline from the
persisted pipeline artifacts --- \emph{no API keys and no internet access are required}
to reproduce the figures, tables, and the interactive app in this report. Secrets live
only in an untracked local \code{.env}; the shared \code{.env.example} is a blank
template. The steps below reproduce the deliverable on a clean machine.

\paragraph{1. Create an isolated environment and install.}
The core install is enough to launch the dashboard and run the figure/test suite; the
bracketed extras add the optional collectors and developer tooling.
{\small\begin{verbatim}
python -m venv .venv
# Windows:        .venv\Scripts\activate
# macOS / Linux:  source .venv/bin/activate
pip install -e .                          # core: dashboard + analysis
pip install -e ".[dev,nlp,telegram,scraping]"   # everything (optional)
\end{verbatim}}

\paragraph{2. Launch the interactive dashboard.}
A single command builds the served payload from the artifacts in \code{data/} and opens
the browser on a local, zero-dependency origin (CORS-free, no external calls).
{\small\begin{verbatim}
python scripts/serve_app.py        # cross-platform
run_app.cmd                        # Windows one-click equivalent
# -> http://127.0.0.1:8000/
\end{verbatim}}

\paragraph{3. Regenerate every figure and table.}
All quantitative results are derived from the persisted artifacts by a single released
step, so nothing in this report is hand-transcribed.
{\small\begin{verbatim}
python scripts/_paper_stats.py     # refresh the stats JSONs from data/
python scripts/_paper_figs.py      # writes pdf_latex/figures/*.pdf
pytest                             # run the test suite (757 passing)
polysignal --help                  # explore the pipeline CLI
\end{verbatim}}

\paragraph{4. (Optional) rebuild the report PDF.}
The \code{pdf\_latex/} folder compiles with any standard TeX toolchain --- \code{pdflatex
+ bibtex}, \code{latexmk}, \code{tectonic}, or Overleaf --- using only standard packages.
{\small\begin{verbatim}
cd pdf_latex && latexmk -pdf main.tex     # or: tectonic main.tex
\end{verbatim}}

\noindent\textbf{What the release bundle contains.} The publishable archive ships the
source (\code{src/}, \code{scripts/}, \code{webapp/}, \code{tests/}), the report
(\code{pdf\_latex/} with \code{main.pdf}), the documentation (\code{docs/}), and the
\emph{minimal} set of \code{data/} artifacts the dashboard reads (timelines, social
pressure, stance labels, accepted evidence, market metadata, post enrichment, and the
analytics/backtest summaries --- about 60\,MB). The multi-gigabyte raw corpus, caches,
virtual environment, and portable tool downloads are deliberately excluded; they are
regenerable and carry no analytical value beyond what the shipped artifacts already
summarise.


\chapter{AI-Assisted Development: Tools and Repositories}
\label{ch:aitools}
In the interest of transparency, this appendix documents the AI-assisted development
workflow behind PolySignal. The \emph{scientific} results are independent of it --- every
number in this report is regenerated deterministically from the persisted artifacts
(\Cref{ch:reproduce}), and no AI agent ever produced a result: the agents wrote
\emph{code}, the code produced the artifacts, and the artifacts produced every figure and
table. What the AI workflow bought was engineering throughput \emph{with} discipline.

\section{AI coding assistants}
Three command-line coding agents were used in deliberately complementary roles, so that
no single model's blind spots went unchecked:
\begin{description}[leftmargin=1.4em,style=nextline]
  \item[Claude Code (Anthropic).] The primary agent for implementation, refactoring, and
  this report. It carried the bulk of the pipeline --- ingestion, contract--post linking,
  outcome-aware stance, the Social Pressure Index, and the leakage-controlled backtest ---
  drove the test-first workflow, and authored and revised the \LaTeX{} of this document.
  \item[Codex (OpenAI).] A second pair of eyes for focused code generation and
  cross-checking: alternative implementations of tricky functions, and independent review
  of edge cases and failure modes.
  \item[Gemini CLI (Google).] Used for large-context sweeps across the codebase and for
  corroborating design decisions, providing a third independent perspective on
  architecture and data-handling choices.
\end{description}
Running three independent agents --- rather than trusting one --- turned model
disagreement into a useful signal: where they agreed we proceeded with confidence; where
they diverged, the divergence flagged exactly the spots that warranted a human decision.

\section{Supporting repositories and tooling}
\begin{description}[leftmargin=1.4em,style=nextline]
  \item[ECC --- a comprehensive ``everything Claude Code'' framework.] ECC is an open
  repository --- the winning entry of a 2026 hackathon --- that packages a complete Claude
  Code setup: agent definitions, multi-agent orchestration rules, reusable skills and
  commands, and a full set of engineering conventions (coding style and immutability, a
  mandatory test-coverage bar, security checks against hard-coded secrets and unvalidated
  input, and a code-review gate before any change is accepted). We adopted ECC to drive
  and standardise the agentic workflow; it is what kept three different agents producing
  consistent, reviewable code rather than three divergent styles, and why the repository
  ships with \textbf{757 passing tests} and a uniform module layout.
  \item[Graphify.] A repository-graph tool (configured through \code{.graphifyignore},
  output under \code{graphify-out/}) that builds a navigable map of the codebase --- its
  modules, symbols, and their relationships. It gave the agents structural context for
  navigating a large repository, so edits respected the existing architecture
  (\Cref{ch:architecture}) instead of reinventing it.
\end{description}

\section{How the tools were leveraged}
The workflow composed these pieces into a single loop: plan a change with one agent,
implement it test-first, cross-review it with a second agent, and keep every change behind
the ECC review-and-coverage gate, using the Graphify map to land edits in the right place.
The result is a large, well-tested, well-documented codebase and a fully reproducible
report, produced quickly but without surrendering the scientific honesty that is the whole
point of a null-result study. The AI tooling accelerated the \emph{engineering}; the
\emph{evidence} remains entirely the data's.
